\documentclass[11pt]{article}
\usepackage[T1]{fontenc}
\usepackage[utf8]{inputenc}
\usepackage{lmodern,amsmath,amssymb,amsthm,microtype}
\usepackage[a4paper,margin=27mm,top=23mm,bottom=25mm]{geometry}
\usepackage{xcolor,hyperref}
\definecolor{linkblue}{HTML}{14576D}
\hypersetup{colorlinks=true,allcolors=linkblue,pdftitle={Sharp double-logarithmic growth of finite restricted orders},pdfauthor={Leo Gao (AI-assisted)}}
\usepackage{fancyhdr}
\pagestyle{fancy}\fancyhf{}\fancyhead[L]{\small\itshape Restricted orders of additive bases}\fancyhead[R]{\small Leo Gao}
\fancyfoot[C]{\thepage}\setlength{\headheight}{14pt}
\newtheorem{theorem}{Theorem}
\newtheorem{lemma}[theorem]{Lemma}
\newtheorem{corollary}[theorem]{Corollary}
\newcommand{\N}{\mathbb N}
\newcommand{\Z}{\mathbb Z}
\newcommand{\lowd}{\underline{d}}
\newcommand{\ord}{\operatorname{ord}}
\newcommand{\ordR}{\operatorname{ord}_{R}}
\newcommand{\set}[1]{\{#1\}}
\DeclareMathOperator{\lcm}{lcm}
\title{Sharp double-logarithmic growth\protect\\ of finite restricted orders}
\author{Leo Gao}
\date{27 November 2026 \enspace$\cdot$\enspace AI-assisted research manuscript}
\begin{document}
\maketitle
\begin{abstract}
Let $F(h)$ be the largest finite restricted order of an additive basis of ordinary order $h$. We prove $F(h)\le \exp(O(2^h h\log(2h)))$. Combined with a previously supplied mixed-radix construction, this gives $\log\log F(h)=h\log 2+O(\log(2h))$. The new step groups sums by how many pairwise disjoint representations they admit, and allocates a different number of summands to each group. This also determines the triple-logarithmic growth rate of the least modulus that tests existence simultaneously for all bases of order at most $h$. The result concerns boundedness and existence, not the classification of equality at a prescribed order.
\end{abstract}
\section{Orders and a robust progression}
For $A\subseteq\N_{>0}$, let $S_h(A)$ and $R_h(A)$ be the sums of at most $h$ terms with repetition allowed and with distinct terms required, respectively. The empty sum is included; adjoining $0$ to $A$ changes neither set. The ordinary and restricted orders, $\ord(A)$ and $\ordR(A)$, are the least $h\ge1$ for which the corresponding sumset is cofinite, or $\infty$ if no such $h$ exists. The four original questions appear in Erdős and Graham \cite[p.~206, remark~3]{EG}. Chen, Yang and Zhao \cite{CYZ} characterize existence via all moduli and settle both finite-deletion questions. We address the size of the largest finite restricted order. All logarithms are natural unless otherwise indicated.

\begin{theorem}[Robust progression]\label{thm:main}
Suppose $\lowd(R_h(A))>\delta$, where $0<\delta<1$ and $\lowd$ denotes lower asymptotic density. There are integers $K,g,N_0\ge1$ and $a\bmod g$ such that, for every $E\subseteq A$ with $|E|\le g-1$,
\[
 [N_0,\infty)\cap(a+g\Z)\ \subseteq\ R_K(A\setminus E),\qquad
 \log(K+g)\ll 2^h\bigl(\log(1/\delta)+\log(2h)\bigr).
\]
The implicit constant is absolute. The onset $N_0$ may depend on $A$, but not on $E$.
\end{theorem}

We use several popularity thresholds per rank to obtain the $2^h$ in the exponent, building on the disjoint-representation method of \cite[Lemmas~3--4]{CYZ}. A preliminary construction in \cite{Gao} supplied the lower bound and an earlier $16^h$ upper estimate. Neither that lower construction nor the underlying density and addition theorems is new here.

\section{One active popularity scale at each length}
Set $D=64h^2$, $L_j=2^j$, $Q_j=D L_j$. Starting from $\alpha_h=\delta$, for $r=h,\ldots,1$ define
\begin{equation}\label{rec}
 J_r=\lceil\log_2(256/\alpha_r)\rceil,\qquad
 \alpha_{r-1}=\frac{\alpha_r^2}{8192\,rD(J_r+1)}.
\end{equation}
Only the finitely many pairs $1\le r\le h$, $0\le j\le J_r$ are used. Let $B_{r,j}$ consist of positive integers admitting $Q_j$ mutually disjoint representations, each a nonempty subset of $A$ with at most $2r$ elements. Write $C_{r,j}=B_{r,j}\cup\{0\}$.

Take $N$ beyond the density onset and large enough that $N\ge\max_{0\le s\le h}1/\alpha_s$. Choose the smallest $r$ with
\[
 M=|R_r(A)\cap[1,N]|\ge\alpha_rN.
\]
Such an $r$ exists by hypothesis. Put $e=|R_{r-1}(A)\cap[0,N]|$. Then $e\le2\alpha_{r-1}N$ by minimality and the possible contribution of $0$. This includes $r=1$, when $e=1$. Select one support $F_x\subseteq A$ of size at most $r$ for each of these $M$ integers $x$. For $1\le z\le 2N$, let $p_z$ be the number of ordered pairs of these integers with $x+y=z$ and $F_x\cap F_y=\varnothing$. We use
\begin{equation}\label{pairs}
 p_z\le M,\qquad \sum_z p_z\ge M^2-rMe,\qquad p_z>4rQ_je\quad\Longrightarrow\quad z\in B_{r,j}.
\end{equation}
Here is the counting argument of \cite[Lemma~3]{CYZ}. Any label lies in at most $e$ supports, since removing it gives different elements of $R_{r-1}(A)\cap[0,N]$. Thus at most $rMe$ ordered pairs have intersecting supports. For a fixed $z\notin B_{r,j}$, a maximal family of disjoint unions $F_x\cup F_y$ of counted pairs has at most $Q_j-1$ members. Its union contains at most $2r(Q_j-1)$ labels and meets every counted pair. Each such label hits at most $2e$ pairs, because $x$ determines $y$ at this fixed $z$, and conversely. This proves \eqref{pairs}.

\begin{lemma}[Adaptive popularity]\label{active}
At each sufficiently large $N$, some permitted pair $(r,j)$ is \emph{active}, that is,
\begin{equation}\label{activeeq}
 L_j |B_{r,j}\cap[1,2N]|\ \ge 16N .
\end{equation}
\end{lemma}
\begin{proof}
Fix $r$ chosen above and abbreviate $J=J_r$. Suppose that no $j\le J$ is active. Write $t_j=4rQ_je$. The entries $p_z\le t_0$ contribute at most $2Nt_0=8rDeN$ in total. For each $j<J$ the band $t_j<p_z\le t_{j+1}$ contributes at most
\[
 t_{j+1}|B_{r,j}\cap[1,2N]|<128rDeN.
\]
Above $t_J$, each $p_z\le M$ and there are fewer than $16N/L_J$ eligible integers; this contribution is at most $M^2/16$ by $L_J\ge256/\alpha_r$. This last bound does not require $t_J\ge M$, so also applies when $e$ is unusually small. Adding the bands gives
\[
 \sum_z p_z\le128rDeN(J+1)+M^2/16
 \le M^2/32+M^2/16
\]
by \eqref{rec}. But $re/M\le 2r\alpha_{r-1}/\alpha_r<1/2$, contradicting \eqref{pairs}.
\end{proof}
The threshold $Q_j$ and the number of summand copies $L_j$ grow together. Thus each band contributes the same error bound, rather than an error proportional to the largest threshold.

\section{Bounding the modulus and avoiding deletions}
Let $j_*$ be the smallest index at which some pair $(r_*,j_*)$ is active for infinitely many $N$, and fix such a pair. All indices below $j_*$ are eventually inactive, since the family is finite. A pair with $B_{r,j}$ finite is likewise eventually inactive. Let $\Gamma$ be the pairs with $j\ge j_*$ and $B_{r,j}$ infinite. By Lemma~\ref{active}, some pair of $\Gamma$ is active at every sufficiently large $N$; it need not be the same pair at different $N$.

Put $S=\sum_{(r,j)\in\Gamma}L_j C_{r,j}$, an ordinary sumset with repetitions allowed. From \eqref{activeeq} and monotonicity between $2N$ and $2N+1$,
\[
 \liminf_{X\to\infty}\frac{1}{X}
    \sum_{(r,j)\in\Gamma} L_j |C_{r,j}\cap[0,X]|\ \ge 8.
\]
All summand sets are infinite and at least two copies occur: an active pair necessarily has $L_j\ge8$. Kneser's asymptotic addition theorem in its mixed counting-function form \cite{Kneser,CYZ} implies that $S$ is eventually periodic, since its own lower density is at most $1$. Explicitly, the comparison is with the liminf of the \emph{sum} of the counting functions, not the sum of their separate liminfs. As each $C_{r,j}$ contains $0$ and is infinite, $S$ contains an eventual progression $a\bmod m$, with $m$ so far unbounded.

Set $L=L_{j_*}$ and let $B$ denote the residues of $C_{r_*,j_*}$ modulo $m$. Across the infinitely many scales where this pair is active, its upper density is at least $8/L$. Thus $|B|\ge8m/L$. The subgroup $H$ generated by $B$ in $\Z/m\Z$ has index $g\le L/8$, and
\[
 LB=H.
\]
Indeed, otherwise let $V$ be the stabilizer of $LB$ in $H$. Since $B$ contains $0$ and generates $H$, it meets $k\ge2$ cosets of $V$. Finite Kneser gives $|LB|\ge(L(k-1)+1)|V|\ge (L/2)|B|\ge4m$, impossible. Choose a sum of $L$ members of $C_{r_*,j_*}$ for each residue in $H$. Adding these finitely many shifts to the eventual progression modulo $m$ fills an eventual progression modulo $g$. Hence this progression lies in
\[
 T=S+L C_{r_*,j_*}.
\]

For $E\subseteq A$ with $|E|\le g-1$, take an ordinary representation of $t\in T$ by the specified sets, ordered by increasing layer $j$. At each rank of layer $i$ there are at most $2L_i$ copies, including the reserve, each to be replaced by at most $2h$ labels. Before a term from layer $j$ is treated, the number of forbidden labels (used labels and $E$) is less than
\[
 g+\sum_{i=j_*}^{j}4h^2L_i<L_j+8h^2L_j<Q_j .
\]
Each positive term has $Q_j$ mutually disjoint supports, one of which therefore avoids all forbidden labels. A zero term uses none. The total cost is less than $K=8h^2L_{J_1}+1$ since $J_1$ is maximal. This proves the desired containment simultaneously over $E$ with a common onset.

\section{Quantitative consequences}
Put $y_r=\log(1/\alpha_r)$. Equation~\eqref{rec} implies with an absolute $C$ that
\[
 y_{r-1}\le2y_r+C\bigl(\log(2h)+\log(2+y_r)\bigr).
\]
First bound $\log(2+y)\le y+O(1)$ to get $y_{h-k}+1\le C_1^k C_2(y_h+\log(2h)+1)$ for absolute $C_1,C_2$. Substituting this preliminary bound inside the logarithmic error in the original recursion, the error at step $k$ is $O(\log(2h)+\log(2+y_h)+k)$. After normalizing the iteration by powers of $2$, these errors sum geometrically because $\sum_{k\ge0}(k+1)2^{-k}<\infty$. Thus $y_1\ll 2^h(y_h+\log(2h))$. Since $L_{J_1}<512\exp(y_1)$ and $g\le L_{J_1}$, this proves Theorem~\ref{thm:main}.

\begin{corollary}\label{growth}
Let $F(h)$ be the maximum finite restricted order among bases of ordinary order exactly $h$. Then
\[
 2^{2^{h-5}+h-4}\le F(h)\le\exp\!\bigl(O(2^h h\log(2h))\bigr)\qquad(h\ge5).
\]
In particular $\displaystyle\lim_{h\to\infty}\frac{\log\log F(h)}h=\log 2$.
\end{corollary}
\begin{proof}
For $\ord(A)\le h$, the density theorem of Hegyvári, Hennecart and Plagne \cite{HHP,Heg} gives
\[
 \lowd(R_h(A))\ge c_h:=h^{-h}\exp(-\pi\sqrt{2h/3}).
\]
To reconcile conventions, apply the exact-$h$ theorem to $(A\cup\{0\})+1$: its ordinary $h$-sumset is cofinite. Translating back by $-h$ and dropping any zero gives a restricted sum of length at most $h$. Use $\delta=c_h/2$ in Theorem~\ref{thm:main}. If $\ordR(A)<\infty$, arbitrary finite subset sums cover every residue modulo $g$. In this quotient a shortest correcting subset for a prescribed residue has at most $g-1$ labels: otherwise two partial sums, including the initial zero, agree. For each target class choose one correction to the progression class. For all sufficiently large targets, subtract it, use the progression outside its support, and append the correction. This proves $\ordR(A)\le K+g-1$ and the stated upper bound. The lower estimate is the earlier mixed-radix construction \cite[Sec.~4]{Gao}. Taking two logarithms yields the limit.
\end{proof}

For fixed $h$, define $U(h)$ to be the least positive integer $u$ such that every basis $A$ of ordinary order at most $h$ has finite restricted order if and only if its unrestricted finite subset sums $\Sigma(A)$ cover all residues modulo $u$.
\begin{corollary}
The modulus $U(h)$ exists, and $\displaystyle\lim_{h\to\infty}h^{-1}\log\log\log U(h)=\log 2$.
\end{corollary}
\begin{proof}
One admissible modulus is $\lcm(1,\ldots,L_{J_1})$ with $\delta=c_h/2$. Indeed any $g$ from the theorem divides it; completeness modulo the larger modulus guarantees corrections modulo $g$. Conversely, a cofinite restricted sumset meets every residue. The bound for $L_{J_1}$ and $\log\lcm(1,\ldots,b)\le b\log b$ give $\log\log\log U(h)\le h\log 2+O(\log(2h))$ once $U(h)$ is sufficiently large.

For the lower bound ($h\ge5$), use the denominations from \cite{Gao}, now with a different tail. Set $d=h-4$, $P_0=1$, $P_i=2^{2^{i-1}+i}$ for $1\le i\le d$, $P=P_d$, $u_i=P_i/P_{i-1}$, $w_i=P/P_i$, and
\[
 C=\{P\}\cup\{k w_i:1\le i\le d,\ 1\le k<u_i\},\qquad S_C=\sum_{c\in C}c .
\]
Then $2P\le S_C<7P/2$: in $S_C/P=1+\frac12\sum_i(u_i-1)/P_{i-1}$ the first place adds $3/2$, and for $i\ge2$ we have $u_i/P_{i-1}=2^{2-i}$, so the rest adds less than $1$. Also $\Sigma(C)=[0,S_C]\cap\Z$. To check the latter assertion, add the denominations from smallest upward. At the start of a place with unit $w_i$, the lower places already cover $[0,w_i-1]$ by their digit expansions. Adding $w_i,2w_i,\ldots,(u_i-1)w_i$ successively cannot create a gap: before adding $k w_i$ the previously covered interval has upper end at least $kw_i-1$. This also applies to the last coin $P$.

Fix $S_C+2\le q\le4P$. Put $A_q=C\cup q\N_{>0}$. Every residue in $[0,q-1]$ uses at most three copies of $P$ and one coin at each digit, so an extra tail term gives $\ord(A_q)\le d+4=h$. But $q-1\notin\Sigma(A_q)\bmod q$. For any proposed universal modulus $M$ with $q\nmid M$, put $v=\gcd(q,M)\le q/2\le2P$. The subset sums of $C$ cover all classes modulo $v$. Every residue in the subgroup $v\Z/M\Z$ is represented by one positive multiple of $q$, choosing a sufficiently large multiple also when the residue is zero. This label is outside $C$ because $q>P$. Thus $\Sigma(A_q)$ is complete modulo $M$, an incorrect test outcome. Every universal modulus is therefore divisible by \emph{every integer} $q$ in the indicated interval. By the prime number theorem the primes in $(7P/2+1,4P]$ alone imply $\log U(h)\gg P$ for all sufficiently large $h$. The claimed lower limit follows.
\end{proof}

\section*{Verification of the earlier lower construction}
We record the argument from \cite[Sec.~4]{Gao} so the use of its lower bound in Corollary~\ref{growth} can be checked independently. Keep the denominations $C$ and period parameter $P$ above. Set $m=8P$ and
\[
 A=\{2c:c\in C\}\cup\{mn,mn+1:n\ge1\}.
\]
Any residue modulo $m$ is $2t+\epsilon$ with $0\le t<4P$ and $\epsilon\in\{0,1\}$. Its even part uses at most $d+3$ doubled denominations by digit expansion, and a sufficiently large tail term in class $\epsilon$ completes the representation. Thus $\ord(A)\le h$.

For the residue $m-1$, the number of residue-one tail summands must be odd, say $2v+1$. At the level of residues, discard all residue-zero terms, retain one $1$, and replace the other $2v$ ones by $v$ copies of the denomination $2$. Subtract the retained $1$ and halve all remaining terms. This gives a representation of some $t\equiv4P-1\pmod{4P}$ in $C$, of length at most the original length minus one. An ordinary representation of $t$ in $C$ needs at least $\lfloor t/P\rfloor+s(t\bmod P)$ coins, where $s$ counts nonzero mixed-radix digits. To see this, normalize from the lowest place upward. At a place of radix $u$, the digits of $b>0$ coins add up to $ku+v'$ with $0\le v'<u$ and $k\le b-1$. Carry the $k$ units to the next place and retain at most one digit at the current place, without increasing the number of coins. The final normalized count is the claimed bound. Here $t\ge4P-1$, and all $d$ digits of $t\bmod P=P-1$ are nonzero. Thus the original length was at least $1+3+d=h$.

In a distinct representation of the same residue, the finite exceptions contribute at most $2S_C<7P$. Their sum plus the number of residue-one tail terms is nonnegative and congruent to $8P-1$ modulo $8P$. At least $8P-1-2S_C\ge P$ such terms are needed. Conversely, for each class choose its representative $k\in\{1,\ldots,m\}$. Every sufficiently large target of this class is a sum of $k$ distinct residue-one tail elements: fix $k-1$ of them and use the growing complement as the last. Thus the restricted order is finite and at least $P=2^{2^{h-5}+h-4}$.

\paragraph{Scope.} The limiting constants above concern boundedness and a modular test for existence. They do not determine $F(h)$ to a constant factor in its logarithm, give a general test for equality of the two orders, or imply the partition and pair-fiber claims in the supplied material. The present proof is AI-assisted and has not undergone external peer review.

\begin{thebibliography}{9}
\bibitem{EG} P. Erdős and R. L. Graham, \href{https://www.renyi.hu/~p_erdos/1980-22.pdf}{On bases with an exact order}, \emph{Acta Arith.} \textbf{37} (1980), 201--207.
\bibitem{CYZ} S.-Q. Chen, Q.-H. Yang and L. Zhao, \href{https://arxiv.org/html/2611.33191v1}{On a problem of Erdős and Graham related to restricted order}, arXiv:2611.33191v1 (2026), Theorem 2 and Lemmas 1, 3--4.
\bibitem{HHP} N. Hegyvári, F. Hennecart and A. Plagne, \href{https://doi.org/10.1515/crll.2003.055}{A proof of two Erdős conjectures on restricted addition and further results}, \emph{J. Reine Angew. Math.} \textbf{560} (2003), 199--220.
\bibitem{Heg} N. Hegyvári, \href{https://real-d.mtak.hu/1038/7/dc_1260_16_doktori_mu.pdf}{Topics in Combinatorial Number Theory}, dissertation, Hungarian Academy of Sciences (2017), Theorem 4.8, p.~63.
\bibitem{Kneser} M. Kneser, \href{https://doi.org/10.1007/BF01174162}{Abschätzung der asymptotischen Dichte von Summenmengen}, \emph{Math. Z.} \textbf{58} (1953), 459--484.
\bibitem{Gao} L. Gao, \href{https://github.com/Leo66Gao/solved-conjectures/tree/main/48-erdos-338-restricted-orders}{How large can the restricted order of an additive basis be?}, AI-assisted research manuscript (27 November 2026), especially Section 4. The present note uses the supplied manuscript version.
\end{thebibliography}
\end{document}
